\documentclass[a4paper,twoside]{article}

\usepackage{epsfig}
\usepackage{subfigure}
\usepackage{calc}
\usepackage{amssymb}
\usepackage{amstext}
\usepackage{amsmath}
\usepackage{amsthm}
\usepackage{multicol}
\usepackage{pslatex}
\usepackage{apalike}
\usepackage{SCITEPRESS}
\usepackage[small]{caption}

\subfigtopskip=0pt
\subfigcapskip=0pt
\subfigbottomskip=0pt

\begin{document}

\title{Homotopy Surface Cutting Using Edges' Sources in Geodesic Distance}

\author{\authorname{Anuwat Dechvijankit,Hiroshi Nagahashi and Kota Aoki}
\affiliation{Department of Computational Intelligence and Systems Science, Tokyo Institute of Technology, Japan}
\email{dechvijankit.a.aa@m.titech.ac.jp}
}

\keywords{geodesic distance, graph cut, homotopy, parameterization}


%The abstract should summarize the contents of the paper and should contain at least 70 and at most 200 words. The text must be set to 9-point font size.
%
%In order to parameterize any three-dimensional surface like a closed boundary one, we need to convert its polygonal mesh into a disk topology surface. For quality of texturing or re-meshing that uses a parameterization technique, it is more effective if the distortion of two-dimensional manifold planar domain map is as small as possible. We introduce a fast way to create seams along with three-dimensional surface for square boundary of a two-dimensional planar domain map. This paper describes a faster and less error technique than the  existing seam cutting methods by simply observing the location of high distortion area in the planar domain. The distortion rate of our proposed method is low as the original one with less time consuming. Also we discuss to find the best corner points in the square planar domain that enable to delivery of the lowest stretch by avoiding boundary constraint if possible.

%Several tools from topology are useful for mesh processing. These tools often operate on the 1-skeleton of a surface, i.e., the graph of edges embedded in the surface. A common task is to find a collection of edges called a cut graph - cutting along these paths turns the surface into a shape which can be flattened into the plane. This kind of flattening is necessary for texture mapping, remeshing, etc. For instance, the video linked to above shows a cut graph (in yellow) on a two-handled torus, which gets flattened to a circular disk.	One way to find a cut graph is to find a set of loops, no two of which are homologous, which cut the surface into a disk when removed. Intuitively, two loops on a surface are homologous if one can be deformed into the other while always keeping it entirely on the surface. For a closed orientable surface with genus g (i.e., a torus with g handles), there are 2g classes of homologically independent loops. A homology basis consists of one loop from each class. Not every homology basis is a cut graph: some homology bases either disconnect the surface or cut it into a punctured sphere. However, a greedy homotopy basis (described below) will cut the surface into a disk.
\abstract{Topology is a property in surfaces that plays a major role in computer graphics. Processing or analysis between two surfaces generally require both of them to be same topology. There are many tools or applications that require disk topology surfaces as input such as parameterization or remeshing. Therefore, we need convert any surfaces to be same as disk. The common procedure is to define edges graph inside surface that be split into two edges and turns surface into topological disk. We call it as homotopy cutting. Problems become more difficult when dealing with high genus surfaces such as torus. Based on a novel method, we present an enhancement method to define cut graph in high-genus surface for homotopy cutting. By using geodesic properties of each edge, we can generate more suitable edges graph than original method while remains performance and stability as original method.}

\onecolumn \maketitle \normalsize \vfill

\section{\uppercase{Introduction}}
\label{sec:introduction}

\noindent Geometry processing is an important research in 3D computer graphics field. Without efficient algorithms, it is very difficult to develop any kinds of advance application for end-users. Some of important applications in 3D computer graphics, such as texture mapping \cite{Bennis:1991:PSF:127719.122744}, normal mapping \cite{Cohen:1998:AS:280814.280832},remeshing \cite{Hormann00quadrilateralremeshing} and parameterization \cite{Tutte:1963,Floater:1997:PSA:248299.248308} require specific topology of input mesh. There are many cases that topological disk surface is specified for further processing.

 Mesh parameterization is defined as a mapping between
a 3D manifold surface and a suitable target domain.
In general, the mesh parameterization is formulated
as a mapping from 3D triangulated surfaces
to a certain 2D planar domain. However, mesh parameterization
requires 3D surfaces to be topologically
equivalent to a disk without any hole, which can
be easily mapped into a 2D domain. Consequently,
closed-boundary surfaces such as 3D models widely
used cannot be parameterized directly. Moreover, surfaces
which have any hole or being non-genus zero also
cannot be parameterized directly too. To resolve the
above problem, we need to convert a 3D surface into
disk topology by cutting seams into the surface to generate
a boundary for parameterization. Parameterization
between these two domains causes distortion errors
such as stretch. Too short/long seam boundary
and random seam without any strategy may give poor
result of distortion easily.

Several applications of three-dimensional surfaces, such as parameterization [21, 50] and
texture mapping [3, 47], require information about the underlying topological structure in
addition to the geometry. In some cases, we wish to simplify the surface topology, to facilitate
algorithms that can be performed only if the surface is a topological disk.
In the texture mapping problem, we wish to nd a continuous and invertible mapping
from the texture, usually a two-dimensional rectangular image, to the surface. Unfortunately,
if the surface is not a topological disk, no such map exists. In such a case, the only feasible
solution is to cut the surface so that it becomes a topological disk. (Haker et al. [28] present
an algorithm for directly texture mapping models with the topology of a sphere, where the
texture is also embedded on a sphere.) Of course, when cutting the surface, one would like to
nd the best possible cut under various considerations. For example, one might want to cut
the surface so that the resulting surface can be texture mapped with minimum distortion
[21, 27, 50]. To our knowledge, all previous approaches for the texture-mapping problem
either rely on cutting heuristics with no quality guarantees [27, 49, 52] or require the user
to cut the surface beforehand [21, 47].


One of the useful applications of parameterization
is re-meshing. There is a famous re-meshing application
called geometry image [1]. Through the parameterization
of any model into the square planar domain,
we can nd the relationship between positions of surface
domain and planar domain. After that, we can
sample grid points on the square planer domain, and
get position on 3D surface. In order to generate such
geometry image with low distortion, the square planar
domain parameterization plays an important role.
Unlikely as the circular boundary domain or natural
boundary domain, the square boundary domain could
be easily caught with the problem of boundary constraint

\section{\uppercase{Manuscript Preparation}}

\noindent We strongly encourage authors to use this document for the
preparation of the camera-ready. Please follow the instructions
closely in order to make the volume look as uniform as possible
\cite{Moore99}.

Please remember that all the papers must be in English and without
orthographic errors.

Do not add any text to the headers (do not set running heads) and
footers, not even page numbers, because text will be added
electronically.

For a best viewing experience the used font must be Times New
Roman, except on special occasions, such as program code
\ref{subsubsec:program_code}.


\subsection{Manuscript Setup}

\noindent The template is composed by a set of 7 files, in the
following 2 groups:\\
\noindent {\bf Group 1.} To format your paper you will need to copy
into your working directory, but NOT edit, the following 4 files:
\begin{verbatim}
  - apalike.bst
  - apalike.sty
  - article.cls
  - scitepress.sty
\end{verbatim}

\noindent {\bf Group 2.} Additionally, you may wish to copy and edit
the following 3 example files:
\begin{verbatim}
  - example.bib
  - example.tex
  - scitepress.eps
\end{verbatim}


\subsection{Page Setup}

The paper size must be set to A4 (210x297 mm). The document
margins must be the following:

\begin{itemize}
    \item Top: 3,3 cm;
    \item Bottom: 4,2 cm;
    \item Left: 2,6 cm;
    \item Right: 2,6 cm.
\end{itemize}

It is advisable to keep all the given values because any text or
material outside the aforementioned margins will not be printed.

\subsection{First Section}

This section must be in one column.

\vfill
\subsubsection{Title and Subtitle}

Use the command \textit{$\backslash$title} and follow the given structure in "example.tex". The title and subtitle must be with initial letters
capitalized (titlecased). If no subtitle is required, please remove the corresponding \textit{$\backslash$subtitle} command. In the title or subtitle, words like "is", "or", "then", etc. should not be capitalized unless they are the first word of the subtitle. No formulas or special characters of any form or language are allowed in the title or subtitle.

\subsubsection{Authors and Affiliations}

Use the command \textit{$\backslash$author} and follow the given structure in "example.tex".

\subsubsection{Keywords}

Use the command \textit{$\backslash$keywords} and follow the given structure in "example.tex". Each paper must have at least one keyword. If more than one is specified, please use a comma as a separator. The sentence must end with a period.

\subsubsection{Abstract}

Use the command \textit{$\backslash$abstract} and follow the given structure in "example.tex".
Each paper must have an abstract up to 200 words. The sentence
must end with a period.

\subsection{Second Section}

Files "example.tex" and "example.bib" show how to create a paper
with a corresponding list of references.

This section must be in two columns.

Each column must be 7,5-centimeter wide with a column spacing
of 0,8-centimeter.

The section text must be set to 10-point.

Section, subsection and sub-subsection first paragraph should not
have the first line indent.

To remove the paragraph indentation (only necessary for the
sections), use the command \textit{$\backslash$noindent} before the
paragraph first word.

If you use other style files (.sty) you MUST include them in the
final manuscript zip file.

\subsubsection{Section Titles}

The heading of a section title should be in all-capitals.

Example: \textit{$\backslash$section\{FIRST TITLE\}}

\vfill
\subsubsection{Subsection Titles}

The heading of a subsection title must be with initial letters
capitalized (titlecased).

Words like "is", "or", "then", etc. should not be capitalized unless
they are the first word of the subsection title.

Example: \textit{$\backslash$subsection\{First Subtitle\}}

\subsubsection{Sub-Subsection Titles}

The heading of a sub subsection title should be with initial letters
capitalized (titlecased).

Words like "is", "or", "then", etc should not be capitalized unless
they are the first word of the sub subsection title.

Example: \textit{$\backslash$subsubsection\{First Subsubtitle\}}

\subsubsection{Tables}

Tables must appear inside the designated margins or they may span
the two columns.

Tables in two columns must be positioned at the top or bottom of the
page within the given margins. To span a table in two columns please add an asterisk (*) to the table \textit{begin} and \textit{end} command.

Example: \textit{$\backslash$begin\{table*\}}

\hspace*{1.5cm}\textit{$\backslash$end\{table*\}}\\

Tables should be centered and should always have a caption
positioned above it. The font size to use is 9-point. No bold or
italic font style should be used.

The final sentence of a caption should end with a period.

\begin{table}[h]
\caption{This caption has one line so it is
centered.}\label{tab:example1} \centering
\begin{tabular}{|c|c|}
  \hline
  Example column 1 & Example column 2 \\
  \hline
  Example text 1 & Example text 2 \\
  \hline
\end{tabular}
\end{table}

\begin{table}[h]
\caption{This caption has more than one line so it has to be
justified.}\label{tab:example2} \centering
\begin{tabular}{|c|c|}
  \hline
  Example column 1 & Example column 2 \\
  \hline
  Example text 1 & Example text 2 \\
  \hline
\end{tabular}
\end{table}

Please note that the word "Table" is spelled out.


\subsubsection{Figures}

Please produce your figures electronically, and integrate them into
your document and zip file.

Check that in line drawings, lines are not interrupted and have a
constant width. Grids and details within the figures must be clearly
readable and may not be written one on top of the other.

Figure resolution should be at least 300 dpi.

Figures must appear inside the designated margins or they may span
the two columns.

Figures in two columns must be positioned at the top or bottom of
the page within the given margins. To span a figure in two columns please add an asterisk (*) to the figure \textit{begin} and \textit{end} command.

Example: \textit{$\backslash$begin\{figure*\}}

\hspace*{1.5cm}\textit{$\backslash$end\{figure*\}}

Figures should be centered and should always have a caption
positioned under it. The font size to use is 9-point. No bold or
italic font style should be used.

\begin{figure}[!h]
  %\vspace{-0.2cm}
  \centering
   {\epsfig{file = SCITEPRESS.eps, width = 5.5cm}}
  \caption{This caption has one line so it is centered.}
  \label{fig:example1}
 \end{figure}

\begin{figure}[!h]
  \vspace{-0.2cm}
  \centering
   {\epsfig{file = SCITEPRESS.eps, width = 5.5cm}}
  \caption{This caption has more than one line so it has to be justified.}
  \label{fig:example2}
  \vspace{-0.1cm}
\end{figure}

The final sentence of a caption should end with a period.



Please note that the word "Figure" is spelled out.

\subsubsection{Equations}

Equations should be placed on a separate line, numbered and
centered.\\The numbers accorded to equations should appear in
consecutive order inside each section or within the contribution,
with the number enclosed in brackets and justified to the right,
starting with the number 1.

Example:

\begin{equation}\label{eq1}
    a=b+c
\end{equation}

\subsubsection{Program Code}\label{subsubsec:program_code}

Program listing or program commands in text should be set in
typewriter form such as Courier New.

Example of a Computer Program in Pascal:

\begin{small}
\begin{verbatim}
 Begin
     Writeln('Hello World!!');
 End.
\end{verbatim}
\end{small}


The text must be aligned to the left and in 9-point type.

\vfill
\subsubsection{Reference Text and Citations}

References and citations should follow the Harvard (Author, date)
System Convention (see the References section in the compiled
manuscript). As example you may consider the citation
\cite{Smith98}. Besides that, all references should be cited in the
text. No numbers with or without brackets should be used to list the
references.

References should be set to 9-point. Citations should be 10-point
font size.

You may check the structure of "example.bib" before constructing the
references.

For more instructions about the references and citations usage
please see the appropriate link at the conference website.

\section{\uppercase{Copyright Form}}

\noindent For the mutual benefit and protection of Authors and
Publishers, it is necessary that Authors provide formal written
Consent to Publish and Transfer of Copyright before publication of
the Book. The signed Consent ensures that the publisher has the
Author's authorization to publish the Contribution.

The copyright form is located on the authors' reserved area.

The printed form should be completed and signed by one author on
behalf of all the other authors, and uploaded through the authors' reserved area. Alternatively, you can send it to the secretariat by e-mail or fax.

\section{\uppercase{Conclusions}}
\label{sec:conclusion}

\noindent Please note that ONLY the files required to compile your paper should be submitted. Previous versions or examples MUST be removed from the compilation directory before submission.

We hope you find the information in this template useful in the preparation of your submission.

\section*{\uppercase{Acknowledgements}}

\noindent If any, should be placed before the references section
without numbering. To do so please use the following command:
\textit{$\backslash$section*\{ACKNOWLEDGEMENTS\}}


\vfill
\bibliographystyle{apalike}
{\small
\bibliography{genus_cutting}}


\section*{\uppercase{Appendix}}

\noindent If any, the appendix should appear directly after the
references without numbering, and not on a new page. To do so please use the following command:
\textit{$\backslash$section*\{APPENDIX\}}

\vfill
\end{document}

